\scanpage{56}
\mspage{5}

est dans $\mathrm{Cat}$, alors $\underline{\mathrm{Ch}}_\infty(X)\unsure{\text{---}|\longrightarrow}X\times X$ ou $\underline{\mathrm{Ch}}(X)\longrightarrow X\times X$ est une $W$-fibration. Dans les deux cas, \deleted{\unsure{je le}} cela (disons) pour $W_\infty$, je le prouve en prouvant que $p$ est à la fois lisse, et propre. Il faudrait donc revoir les notions de lissité et de propreté dans le contexte de la donnée d'un $W$.

Si je me rappelle bien, la propriété a$'$) était quelque chose d'assez aisément vérifiable, il suffisait de supposer que les
\[
X_y\longrightarrow X_{/y}\quad\text{et }\deleted{\unsure{le}}X_{/y}\mathrel{\text{---}}X_{/y'}
\]
(pour $y\in\mathrm{Ob}\,Y$ resp. $u:y\longrightarrow y'$ dans $\mathrm{Fl}\,Y$) sont dans $W$, et que ça reste vrai par tout chgt de base -- ce qui revient \deleted{\unsure{pratiquement}} à la \deleted{\unsure{axiome}} condition $a_1$ sur $f$.

On pourrait donc formuler une \deleted{\unsure{axi}} condition sur $W$ ainsi :

% slanted
\margin{$\unsure{W(4)|W'(4)}$ Axiome des $W$-fibrations (Voir variante suffisante, dans le NB p. 10) Voir \unsure{cascade|série} de variantes \deleted{\unsure{pages}} 13,14 ci-dessous}

\textcircled{4} \inserted{(faible)} Soit $f:X\longrightarrow Y$ dans $\mathrm{Cat}$ satisfaisant la condition de $W$-fibration $a_1)$, ci-dessus (p. 4), alors elle satisfait même : a) i.e le changement de base par $f$ transforme flèches dans $W$ en flèches dans $W$. \inserted{i.e. une $W$-fibration à gauche est une fibr. ultraforte}

Grâce à la prop. 2, il s'ensuit que la condition duale $a_2)$ sur $f$ implique la condition de $W$-fibration forte, donc l'axiome 4, en apparence asymétrique, \deleted{\unsure{équi}} signifie \unsure{aussi|ainsi} que les \inserted{quatre} conditions de fibration sur $f$ sont équivalentes, il est donc en fait symétrique. \deleted{\unsure{Mai}}

\uline{NB} Je ne vois pas, sans axiomes supplémentaires sur $W$, comment passer de $a_1$ à $a'$ seulement. La difficulté, c'est de montrer, moyennant $a_1$ \unsure{seulement}, que \deleted{\unsure{que}} si $Y$ est $W$-asphérique, l'inclusion $X_y\longrightarrow X$ (où $y\in\mathrm{Ob}\,Y$) est dans $W$. Dans le contexte des dérivateurs $\mathbb{D}$, cela se voyait par voie $\mathbb{D}$-cohomologique.

Cet axiome très fort suffit à impliquer que si $X,Y$ sont $W$-fibrés sur $S$, et si $f:X\longrightarrow Y$ est un $S$-morphisme, qui est dans $W$, alors $f_s:X_s\longrightarrow Y_s$ est dans $W$ $\forall s\in S$ ? Montrons déjà au moins l'inverse.

% archive sheet number: 56
